\begin{itemize}
    \item Se diseñó una Arquitectura Dirigida por Eventos (EDA, por sus siglas en inglés) para integrar las escenas, interacciones y sistemas del videoclip musical en una experiencia de realidad virtual (VR por sus siglas en inglés). La arquitectura distribuyó la comunicación entre productores y consumidores, y usó una configuración compartida y un reloj musical para coordinar la secuencia.
    \item Se desarrolló el sistema de secuenciación temporal y se obtuvo evidencia de su funcionamiento en una ejecución completa. Los registros permitieron seguir los 11 pasos definidos para el flujo, desde la solicitud en el menú hasta el cierre de la experiencia y el retorno al menú. También se correlacionaron 1002 publicaciones y recepciones, resumidas en 19 relaciones entre productores y consumidores. 
    \item Se registró la activación de los seis segmentos y se compararon sus tiempos musicales con las fronteras previstas en la secuencia. Las seis tuvieron datos comparables y la diferencia máxima fue de 17.77 ms. Se contaron 488 publicaciones de \emph{BeatReached}, recibidas por cuatro consumidores y una omisión por \emph{exceeded\_max\_lateness}.
    \item Se evaluó el rendimiento en tres corridas de PC por Meta Horizon Link y tres corridas en la aplicación instalada en Meta Quest 3S. En PC, las medianas de procesamiento de CPU y GPU se mantuvieron entre 1.252 y 1.416 ms, y entre 1.597 y 1.622 ms; también, se registraron descartes de fotogramas. En Quest 3S la mediana fue de 72 FPS en las tres corridas, aunque una parte de la muestra quedo por debajo de ese valor y se observaron algunos picos de GPU. Por lo tanto el funcionamiento de la secuencia y de la arquitectura no implicó un rendimiento uniforme en todas las muestras. 
    \item Se caracterizó la percepción de presencia por medio de I-group Presence Questionnaire (IPQ por sus siglas en inglés). La presencia general obtuvo una media de 2.05 sobre 3 y la presencia espacial una media de 1.97. Y las medidas de involucramiento y realismo de 0.38 y 0.15. En estas muestras las puntuaciones son más altas en la sensación de estar dentro del entorno para que el involucramiento y el realismo percibido. 
    \item El trabajo integró una secuencia musical interactiva, comunicación por eventos, activación temporal de escenas, evaluación de rendimiento y de la percepción de presencia. Los resultados permiten describir cada aspecto, pero no establecer que la arquitectura, el rendimiento o una característica en particular de la experiencia haya causado las puntuaciones de IPQ.
\end{itemize}